\chapter{Weak Solutions, Part I}

In this chapter and the next, we discuss the regularity of weak solutions of
elliptic equations of divergence form. In order to explain ideas clearly, we will
discuss equations of the following form only
\[
-D_j(a_{ij}(x)D_i u)+c(x)u=f(x).
\]

We assume $\Omega$ is a domain in $\R^n$. The function $u\in H^1(\Omega)$ is a weak solution if
it satisfies
\[
\int_\Omega \bigl(a_{ij}D_i u D_j\varphi + c u\varphi\bigr)=\int_\Omega f\varphi
\qquad \text{for any }\varphi\in H^1_0(\Omega),
\]
where we assume

(i) the leading coefficients $a_{ij}\in L^\infty(\Omega)$ are uniformly elliptic, i.e., for some
positive constant $\lambda$ there holds
\[
a_{ij}(x)\xi_i\xi_j\ge \lambda|\xi|^2 \qquad \text{for any }x\in\Omega\text{ and }\xi\in\R^n;
\]

(ii) the coefficient $c\in L^{\frac n2}(\Omega)$ and the nonhomogeneous term $f\in L^{\frac{2n}{n+2}}(\Omega)$.

By the Sobolev embedding theorem, (ii) is the least assumption on $c$ and $f$ to have
a meaningful equation.

We will prove various interior regularity results concerning the solution $u$ if we
have better assumptions on coefficients $a_{ij}$ and $c$ and on the nonhomogeneous term
$f$. Basically there are two class of regularity results, perturbation results and non-
perturbation results. The first is based on regularity assumptions on the leading
coefficients $a_{ij}$, which are assumed to be at least continuous. Under such assumptions, we compare solutions of the underlying equations with harmonic functions,
or solutions of constant coefficient equations. Then, the regularity of solutions
depends on how close they are to harmonic functions or how close the leading
coefficients $a_{ij}$ are to constant coefficients. In this direction, we have Schauder
estimates and $W^{2,p}$ estimates. In this chapter, we only discuss the Schauder es-
timates. For the second kind of regularity, there is no continuity assumption on
the leading coefficients $a_{ij}$. Hence the result is not based on perturbations. The
iteration methods introduced by DeGiorgi and Moser are successful in dealing with
the non-perturbation situation. Results proved by them are fundamental for the
discussion of quasilinear equations, where coefficients depend on solutions. In fact,
the linearity has no bearing in their arguments. This permits an extension of these
results to quasilinear equations with appropriate structure conditions.

Boundary regularities can be discussed in a similar way. We leave details to
readers.

\section*{3.1. Growth of Local Integrals}

Let $B_R(x_0)$ be the ball in $\R^n$ of radius $R$ centered at $x_0$. The well-known
Sobolev theorem asserts that, if $u\in W^{1,p}(B_R(x_0))$ with $p>n$, then $u$ is H\"older
continuous with exponent $\alpha=1-n/p$.

In the first part of this section, we prove a general result, due to S. Campanato,
which characterizes H\"older continuous functions by the growth of their local inte-
grals. This result is very useful for studying the regularity of solutions of elliptic
differential equations. In the second part of this section we prove a result, due
to John and Nirenberg, which gives an equivalent definition of functions of the
bounded mean oscillation.

Let $\Omega$ be a bounded connected open set in $\R^n$ and let $u\in L^1(\Omega)$. For any ball
$B_r(x_0)\subset \Omega$, define
\[
u_{x_0,r}=\frac{1}{|B_r(x_0)|}\int_{B_r(x_0)}u.
\]

We first prove Campanato's characterization of H\"older continuity.

\medskip

\begin{theorem}\label{hanlin:weaksol-part1-campanato}\dcref{Theorem 3.1}\group{ch03-weak-solutions-part-i}\source{han-lin-lecture-notes:page-0058,han-lin-lecture-notes:page-0059}
\textit{Suppose $u\in L^2(\Omega)$ satisfies for some $\alpha\in(0,1)$}
\[
\int_{B_r(x)}|u-u_{x,r}|^2\le M^2r^{n+2\alpha}\qquad \textit{for any }B_r(x)\subset \Omega.
\]
\textit{Then $u\in C^\alpha(\Omega)$, and for any $\Omega'\Subset\Omega$}
\[
\sup_{\Omega'}|u|+\sup_{x,y\in\Omega',\,x\ne y}\frac{|u(x)-u(y)|}{|x-y|^\alpha}\le c(M+\|u\|_{L^2(\Omega)}),
\]
\textit{where $c$ is a positive constant depending only on $n,\alpha,\Omega$ and $\Omega'$.}
\end{theorem}
\begin{proof}
Denote $R_0=\operatorname{dist}(\Omega',\partial\Omega)$. For any $x_0\in \Omega'$ and $0<r_1<r_2\le R_0$, we
have
\[
|u_{x_0,r_1}-u_{x_0,r_2}|^2\le 2\bigl(\,|u(x)-u_{x_0,r_1}|^2+|u(x)-u_{x_0,r_2}|^2\bigr).
\]
By integrating with respect to $x$ in $B_{r_1}(x_0)$, we obtain
\[
|u_{x_0,r_1}-u_{x_0,r_2}|^2\le \frac{2}{\omega_n r_1^n}\left(\int_{B_{r_1}(x_0)}|u-u_{x_0,r_1}|^2+\int_{B_{r_2}(x_0)}|u-u_{x_0,r_2}|^2\right),
\]
and hence
\[
(1)\qquad |u_{x_0,r_1}-u_{x_0,r_2}|^2\le c(n)M^2r_1^{-n}\bigl(r_1^{n+2\alpha}+r_2^{n+2\alpha}\bigr).
\]
For any $R\le R_0$, with $r_1=R/2^{i+1}$, $r_2=R/2^i$, we get
\[
|u_{x_0,2^{-(i+1)}R}-u_{x_0,2^{-i}R}|\le c(n)2^{-(i+1)\alpha}MR^\alpha,
\]
and therefore for $h<k$
\[
|u_{x_0,2^{-h}R}-u_{x_0,2^{-k}R}|\le \frac{c(n)}{2^{(h+1)\alpha}}MR^\alpha\sum_{i=0}^{k-h-1}\frac{1}{2^{i\alpha}}\le \frac{c(n,\alpha)}{2^{h\alpha}}MR^\alpha.
\]
This shows that $\{u_{x_0,2^{-i}R}\}\subset \R$ is a Cauchy sequence, hence a convergent one.
Its limit $\hat{u}(x_0)$ is independent of the choice of $R$, since (1) can be applied with
$r_1=2^{-i}R$ and $r_2=2^{-i}\bar{R}$ whenever $0<R<\bar{R}\le R_0$. Thus we get
\[
\hat{u}(x_0)=\lim_{r\to 0}u_{x_0,r},
\]
with
\[
(2)\qquad |u_{x,r}-\hat{u}(x_0)|\le c(n,\alpha)Mr^\alpha\qquad\text{for any }0<r\le R_0.
\]
Note that $\{u_{x,r}\}$ converges, as $r\to 0+$, in $L^1(\Omega)$ to the function $u$, by the Lebesgue theorem. Then we have $u=\hat{u}$ a.e., and (2) implies that $\{u_{x,r}\}$ converges uniformly to $u(x)$ in $\Omega'$. Since $x\mapsto u_{x,r}$ is continuous for any $r>0$, $u(x)$ is continuous. By (2) we get
\[
|u(x)|\le CMR^\alpha+|u_{x,R}|,
\]
for any $x\in\Omega'$ and $R\le R_0$. Hence $u$ is bounded in $\Omega'$ with
\[
\sup_{\Omega'}|u|\le c(MR_0^\alpha+\|u\|_{L^2(\Omega)}).
\]
Now we prove that $u$ is Hölder continuous. Let $x,y\in\Omega'$ with $R=|x-y|<R_0/2$. Then we have
\[
|u(x)-u(y)|\le |u(x)-u_{x,2R}|+|u(y)-u_{y,2R}|+|u_{x,2R}-u_{y,2R}|.
\]
The first two terms on the right hand side are already estimated in (2). For the last term we write
\[
|u_{x,2R}-u_{y,2R}|\le |u_{x,2R}-u(\zeta)|+|u_{y,2R}-u(\zeta)|.
\]
Integrating with respect to $\zeta$ over $B_{2R}(x)\cap B_{2R}(y)$, which contains $B_R(x)$, yields
\[
|u_{x,2R}-u_{y,2R}|^2\le \frac{2}{|B_R(x)|}\left(\int_{B_{2R}(x)}|u-u_{x,2R}|^2+\int_{B_{2R}(y)}|u-u_{y,2R}|^2\right)
\le c(n,\alpha)M^2R^{2\alpha}.
\]
Therefore we have
\[
|u(x)-u(y)|\le c(n,\alpha)M|x-y|^\alpha.
\]
For $|x-y|>R_0/2$ we obtain
\[
|u(x)-u(y)|\le 2\sup_{\Omega'}|u|\le c\left(M+\frac{1}{R_0^\alpha}\|u\|_{L^2}\right)|x-y|^\alpha.
\]
This finishes the proof. \hfill$\square$
\end{proof}

The Sobolev theorem is an easy consequence of Theorem 3.1. In fact, we have the following result due to Morrey.

\medskip

\begin{corollary}\label{hanlin:weaksol-part1-morrey}\dcref{Corollary 3.2}\group{ch03-weak-solutions-part-i}\source{han-lin-lecture-notes:page-0059}
\textit{Suppose $u\in H^1_{\mathrm{loc}}(\Omega)$ satisfies for some $\alpha\in(0,1)$}
\[
\int_{B_r(x)}|Du|^2\le M^2r^{n-2+2\alpha}\qquad\text{for any }B_r(x)\subset\Omega.
\]
\textit{Then $u\in C^\alpha(\Omega)$, and for any $\Omega'\Subset\Omega$}
\[
\sup_{\Omega'}|u|+\sup_{x,y\in\Omega',\,x\ne y}\frac{|u(x)-u(y)|}{|x-y|^\alpha}\le c(M+\|u\|_{L^2(\Omega)}),
\]
\textit{where $c$ is a positive constant depending only on $n,\alpha,\Omega$ and $\Omega'$.}
\end{corollary}
\begin{proof}
By the Poincar\'e inequality, we obtain
\[
\int_{B_r(x)}|u-u_{x,r}|^2\le c(n)r^2\int_{B_r(x)}|Du|^2\le c(n)M^2r^{n+2\alpha}.
\]
With Theorem 3.1, we have the desired result. \hfill$\square$
\end{proof}

\begin{lemma}\label{hanlin:weaksol-part1-lemma-3-3}\dcref{Lemma 3.3}\group{ch03-weak-solutions-part-i}\source{han-lin-lecture-notes:page-0060}
Suppose $u \in H^1(\Omega)$ satisfies for some $\mu \in [0,n]$
\[
\int_{B_r(x_0)} |Du|^2 \le M r^\mu
\qquad \text{for any } B_r(x_0) \subset \Omega.
\]
Then for any $\Omega' \Subset \Omega$ and any $B_r(x_0) \subset \Omega$ with $x_0 \in \Omega'$
\[
\int_{B_r(x_0)} |u|^2 \le c\Bigl(M + \int_\Omega u^2\Bigr) r^\lambda,
\]
where $\lambda = \mu + 2$ if $\mu < n - 2$ and $\lambda$ is any number in $[0,n]$ if $n - 2 \le \mu < n$, and $c$
is a positive constant depending only on $n$, $\lambda$, $\mu$, $\Omega$ and $\Omega'$.
\end{lemma}
\begin{proof}
As before denote $R_0 = \operatorname{dist}(\Omega', \partial\Omega)$. For any $x_0 \in \Omega'$ and $0 < r \le R_0$,
the Poincar\'e inequality yields
\[
\int_{B_r(x_0)} |u - u_{x_0,r}|^2 \le c r^2 \int_{B_r(x_0)} |Du|^2\,dx \le c(n) M r^{\mu+2}.
\]
This implies
\[
\int_{B_r(x_0)} |u - u_{x_0,r}|^2 \le c(n) M r^\lambda,
\]
where $\lambda$ is as in Lemma 3.3. For any $0 < \rho < r \le R_0$, we have
\[
\int_{B_\rho(x_0)} u^2 \le 2 \int_{B_\rho(x_0)} |u_{x_0,r}|^2 + 2 \int_{B_\rho(x_0)} |u - u_{x_0,r}|^2
\]
\[
\le c(n)\rho^n |u_{x_0,r}|^2 + 2 \int_{B_r(x_0)} |u - u_{x_0,r}|^2
\]
\[
\le c(n)\Bigl(\frac{\rho}{r}\Bigr)^n \int_{B_r(x_0)} u^2 + M r^\lambda,
\]
where we used
\[
|u_{x_0,r}|^2 \le \frac{c(n)}{r^n} \int_{B_r(x_0)} u^2.
\]
Hence we have
\[
(1)\qquad \int_{B_\rho(x_0)} u^2 \le c\Bigl(\frac{\rho}{r}\Bigr)^n \int_{B_r(x_0)} u^2 + M r^\lambda
\qquad \text{for any } 0 < \rho < r \le R_0.
\]
We note $\lambda \in (0,n)$. Now we claim
\[
(2)\qquad \int_{B_\rho(x_0)} u^2 \le c\Bigl(\frac{\rho}{r}\Bigr)^\lambda \int_{B_r(x_0)} u^2 + M \rho^\lambda
\qquad \text{for any } 0 < \rho < r \le R_0.
\]
By choosing $r = R_0$, we obtain
\[
\int_{B_\rho(x_0)} u^2 \le c\rho^\lambda\Bigl(\int_\Omega u^2 + M\Bigr)
\qquad \text{for any } \rho \le R_0.
\]
In order to get (2) from (1), we apply the following technical lemma to the function
$\phi(r)=\int_{B_r(x_0)} u^2$.
\end{proof}

\begin{lemma}\label{hanlin:weaksol-part1-lemma-3-4}\dcref{Lemma 3.4}\group{ch03-weak-solutions-part-i}\source{han-lin-lecture-notes:page-0061}
Let $\phi(t)$ be a nonnegative and nondecreasing function on $[0,R]$.
Suppose for some nonnegative constants $A,B,\alpha,\beta$ with $\beta < \alpha$
\[
\phi(\rho) \leq A\left(\left(\frac{\rho}{r}\right)^{\alpha}+\epsilon\right)\phi(r) + Br^{\beta}
\qquad \text{for any } 0<\rho \leq r \leq R.
\]
Then for any $\gamma \in (\beta,\alpha)$, there exists a constant $\epsilon_0$ such that, if $\epsilon < \epsilon_0$,
\[
\phi(\rho) \leq c\left(\left(\frac{\rho}{r}\right)^{\gamma}\phi(r) + Br^{\beta}\right)
\qquad \text{for any } 0<\rho \leq r \leq R,
\]
where $\epsilon_0$ and $c$ are positive constants depending only on $A,\alpha,\beta$ and $\gamma$. In particular,
\[
\phi(r) \leq c\left(\frac{\phi(R)}{R^{\gamma}}r^{\gamma} + Br^{\beta}\right)
\qquad \text{for any } 0<r \leq R.
\]
\end{lemma}
\begin{proof}
For any $\tau \in (0,1)$ and $r < R$, we have
\[
\phi(\tau r) \leq A\tau^{\alpha}(1+\epsilon\tau^{-\alpha})\phi(r) + Br^{\beta}.
\]
Choose $\tau < 1$ such that $2A\tau^{\alpha} = \tau^{\gamma}$ and assume $\epsilon_0\tau^{-\alpha} < 1$. Then we get for every
$r < R$
\[
\phi(\tau r) \leq \tau^{\gamma}\phi(r) + Br^{\beta},
\]
and therefore for any integer $k > 0$
\[
\phi(\tau^{k+1}r) \leq \tau^{\gamma}\phi(\tau^k r) + B\tau^{k\beta}r^{\beta}
\]
\[
\leq \tau^{(k+1)\gamma}\phi(r) + B\tau^{k\beta}r^{\beta}\sum_{j=0}^{k}\tau^{j(\gamma-\beta)}
\]
\[
\leq \tau^{(k+1)\gamma}\phi(r) + \frac{B\tau^{k\beta}r^{\beta}}{1-\tau^{\gamma-\beta}}.
\]
Choosing $k$ such that $\tau^{k+2}r < \rho \leq \tau^{k+1}r$, we obtain
\[
\phi(\rho) \leq \frac{1}{\tau^{\gamma}}\left(\frac{\rho}{r}\right)^{\gamma}\phi(r) +
\frac{B\rho^{\beta}}{\tau^{2\beta}(1-\tau^{\gamma-\beta})}.
\]
This finishes the proof.

In the rest of this section, we discuss functions of the bounded mean oscillation
(BMO). The following result is due to John and Nirenberg and referred to as the
John-Nirenberg lemma.
\end{proof}

\begin{theorem}\label{hanlin:weaksol-part1-john-nirenberg}\dcref{Theorem 3.5}\group{ch03-weak-solutions-part-i}\source{han-lin-lecture-notes:page-0061,han-lin-lecture-notes:page-0062,han-lin-lecture-notes:page-0063}
\textit{Suppose $u \in L^{1}(\Omega)$ satisfies}
\[
\int_{B_r(x)} \lvert u-u_{x,r}\rvert \leq Mr^n
\qquad \text{for any } B_r(x)\subset \Omega.
\]
\textit{Then for any } $B_r(x)\subset \Omega$
\[
\int_{B_r(x)} e^{\frac{p_0}{M}\lvert u-u_{x,r}\rvert} \leq Cr^n,
\]
\textit{where } $p_0$ \textit{and} $C$ \textit{are positive constants depending only on} $n$.
\end{theorem}
\begin{remark}\label{hanlin:weaksol-part1-bmo}\dcref{Remark 3.6}\group{ch03-weak-solutions-part-i}\source{han-lin-lecture-notes:page-0061}
Functions satisfying the condition of Theorem 3.5 are called
functions of the bounded mean oscillation (BMO). We note that
\[
L^{\infty} \text{ is a proper subset of } \BMO.
\]
The function $u(x)=\log(x)$ in $(0,1)\subset \mathbb{R}$ is in BMO but not in $L^{\infty}$.
\end{remark}

For convenience, we use cubes instead of balls. We need the Calderon-Zygmund decomposition in the proof of Theorem 3.5. First, we introduce some terminology.
Take the unit cube $Q_0$. Cut it equally into $2^n$ cubes, which we take as the first generation. Do the same cutting for these small cubes to get the second generation. Continue this process. These cubes (from all generations) are called dyadic cubes.
Any $(k+1)$-generation cube $Q$ arises from some $k$-generation cube $\widetilde{Q}$, which is called the predecessor of $Q$.

\begin{lemma}\label{hanlin:weaksol-part1-cz}\dcref{Lemma 3.7}\group{ch03-weak-solutions-part-i}\source{han-lin-lecture-notes:page-0062}
\textit{Suppose $f \in L^1(Q_0)$ is nonnegative and $\alpha > |Q_0|^{-1}\int_{Q_0} f$ is a fixed constant. Then there exists a sequence of (nonoverlapping) dyadic cubes $\{Q_j\}$ in $Q_0$ such that}
\[
f(x) \leq \alpha \quad \text{a.e. in } Q_0 \setminus \bigcup_j Q_j,
\]
\textit{and}
\[
\alpha \leq \frac{1}{|Q_j|}\int_{Q_j} f\,dx < 2^n\alpha.
\]
\end{lemma}
\begin{proof}
Cut $Q_0$ into $2^n$ dyadic cubes and keep the cube $Q$ if $|Q|^{-1}\int_Q f \geq \alpha$.
Continue cutting for others, and always keep the cube $Q$ if $|Q|^{-1}\int_Q f \geq \alpha$ and cut the rest. Let $\{Q_j\}$ be the cubes we have kept during this infinite process. We only need to verify
\[
f(x) \leq \alpha \quad \text{a.e. in } Q_0 \setminus \bigcup_j Q_j.
\]
Let $F = Q_0 \setminus \bigcup_j Q_j$. For any $x \in F$, from the way we collect $\{Q_j\}$, there exists a sequence of cubes $Q^i$ containing $x$ such that
\[
\frac{1}{|Q^i|}\int_{Q^i} f < \alpha,
\]
and
\[
\diam(Q^i) \to 0 \text{ as } i \to \infty.
\]
The Lebesgue density theorem implies
\[
f \leq \alpha \quad \text{a.e. in } F.
\]
This finishes the proof.
\end{proof}

\begin{proof}[Proof of Theorem 3.5]
Assume $\Omega = Q_0$. We may rewrite the assumption in terms of cubes as follows
\[
\int_Q |u-u_Q| < M|Q| \qquad \text{for any } Q \subset Q_0.
\]
We prove that there exist two positive constants $c_1$ and $c_2$, depending only on $n$, such that for any $Q \subset Q_0$
\[
|\{x \in Q; |u-u_Q| > t\}| \leq c_1|Q|\exp\!\left(-\frac{c_2}{M}t\right).
\]
Then Theorem 3.5 follows easily.
Without loss of generality, we assume $M = 1$. Choose
\[
\alpha > 1 \geq |Q_0|^{-1}\int_{Q_0} |u-u_{Q_0}|\,dx.
\]
Apply the Calderon-Zygmund decomposition to $f = |u - u_{Q_0}|$. There exists a
sequence of (nonoverlapping) cubes $\{Q_j^{(1)}\}_{j=1}^{\infty}$ such that
\[
\alpha \leq \frac{1}{|Q_j^{(1)}|}\int_{Q_j^{(1)}} |u - u_{Q_0}| < 2^n\alpha,
\]
and
\[
|u(x) - u_{Q_0}| \leq \alpha \quad \text{a.e. } x \in Q_0 \setminus \bigcup_{j=1}^{\infty} Q_j^{(1)}.
\]
This implies
\[
\sum_j |Q_j^{(1)}| \leq \frac{1}{\alpha}\int_{Q_0} |u - u_{Q_0}| \leq \frac{1}{\alpha}|Q_0|,
\]
and
\[
|u_{Q_j^{(1)}} - u_{Q_0}| \leq \frac{1}{|Q_j^{(1)}|}\int_{Q_j^{(1)}} |u - u_{Q_0}|\,dx \leq 2^n\alpha.
\]
The definition of the BMO norm implies for each $j$
\[
\frac{1}{|Q_j^{(1)}|}\int_{Q_j^{(1)}} |u - u_{Q_j^{(1)}}|\,dx \leq 1 < \alpha.
\]
Apply the decomposition procedure above to $f = |u - u_{Q_j^{(1)}}|$ in
$Q_j^{(1)}$. There exists a sequence of (nonoverlapping) cubes
\[
\{Q_j^{(2)}\} \text{ in } \bigcup_j Q_j^{(1)} \text{ such that}
\]
\[
\sum_j |Q_j^{(2)}| \leq \frac{1}{\alpha}\sum_j \int_{Q_j^{(1)}} |u - u_{Q_j^{(1)}}|
\leq \frac{1}{\alpha}\sum_j |Q_j^{(1)}|
\leq \frac{1}{\alpha^2}|Q_0|,
\]
and
\[
|u(x) - u_{Q_j^{(1)}}| \leq \alpha \quad \text{for a.e. } x \in Q_j^{(1)} \setminus \bigcup_j Q_j^{(2)},
\]
which implies
\[
|u(x) - u_{Q_0}| \leq 2\cdot 2^n\alpha \quad \text{for a.e. } x \in Q_0 \setminus \bigcup_j Q_j^{(2)}.
\]
Continue this process. For any integer $k \geq 1$, there exists a sequence of disjoint
cubes $\{Q_j^{(k)}\}$ such that
\[
\sum_j |Q_j^{(k)}| \leq \frac{1}{\alpha^k}|Q_0|,
\]
and
\[
|u(x) - u_{Q_0}| \leq k2^n\alpha \quad \text{for a.e. } x \in Q_0 \setminus \bigcup_j Q_j^{(k)}.
\]
Thus, we obtain
\[
|\{x \in Q_0; |u - u_{Q_0}| > 2^n k\alpha\}| \leq \sum_j |Q_j^{(k)}| \leq \frac{1}{\alpha^k}|Q_0|.
\]
For any $t$, there exists an integer $k$ such that $t \in [2^n k\alpha, 2^n(k+1)\alpha)$.
This implies
\[
\alpha^{-k} = \alpha \alpha^{-(k+1)} = \alpha e^{-(k+1)\log \alpha} \leq \alpha \exp\!\left(-\frac{\log \alpha}{2^n\alpha}t\right).
\]
This finishes the proof.
\end{proof}

\section*{3.2. Hölder Continuity of Solutions}

In this section, we prove the Hölder regularity for solutions. The basic idea is
to freeze the leading coefficients and then to compare solutions with harmonic func-
tions. The regularity of solutions depends on how close solutions are to harmonic
functions. Hence, we need some regularity assumption on the leading coefficients.

Suppose $a_{ij} \in L^\infty(B_1)$ is uniformly elliptic in $B_1$, i.e.,
\[
\lambda |\xi|^2 \le a_{ij}(x)\xi_i\xi_j \le \Lambda |\xi|^2
\qquad \text{for any } x \in B_1, \xi \in \R^n.
\]

In the following, we assume that $a_{ij}$ is at least continuous and that $u \in H^1(B_1)$
satisfies
\[
(\ast)\qquad \int_{B_1} a_{ij}D_i u D_j\varphi + cu\varphi
= \int_{B_1} f\varphi
\qquad \text{for any } \varphi \in H^1_0(B_1).
\]

The main theorem in this section is the following Hölder estimates for solutions.

\begin{theorem}\label{hanlin:weaksol-part1-holder-sol}\dcref{Theorem 3.8}\group{ch03-weak-solutions-part-i}\source{han-lin-lecture-notes:page-0064,han-lin-lecture-notes:page-0065,han-lin-lecture-notes:page-0066,han-lin-lecture-notes:page-0067,han-lin-lecture-notes:page-0068,han-lin-lecture-notes:page-0069,han-lin-lecture-notes:page-0070}
Let $u \in H^1(B_1)$ solve $(\ast)$. Assume $a_{ij} \in C^0(\overline{B_1})$, $c \in L^n(B_1)$ and $f \in L^q(B_1)$ for some $q \in (n/2, n)$. Then $u \in C^\alpha(B_1)$ with $\alpha = 2-n/q \in (0,1)$.

Moreover, there exists an $R_0$ such that, for any $x \in B_{1/2}$ and $r \le R_0$,
\[
\int_{B_r(x)} |Du|^2 \le Cr^{n-2+2\alpha}\bigl(\|f\|_{L^q(B_1)}^2 + \|u\|_{H^1(B_1)}^2\bigr),
\]
where $R_0$ and $C$ are positive constants depending only on $\lambda,\Lambda,\tau$ and $\|c\|_{L^n}$, with $\tau$
denoting the modulus of the continuity of $a_{ij}$, i.e.,
\[
|a_{ij}(x)-a_{ij}(y)| \le \tau(|x-y|)
\qquad \text{for any } x,y \in B_1.
\]

If $c \equiv 0$, we may replace $\|u\|_{H^1(B_1)}$ with $\|Du\|_{L^2(B_1)}$.
\end{theorem}

The idea of the proof is to compare the solution $u$ with harmonic functions
and use the perturbation argument. First, we prove a basic estimate for harmonic
functions.

\begin{lemma}\label{hanlin:weaksol-part1-harmonic-lemma}\dcref{Lemma 3.9}\group{ch03-weak-solutions-part-i}\source{han-lin-lecture-notes:page-0064}
Suppose $\{a_{ij}\}$ is a constant positive definite matrix satisfying for
some $0<\lambda\le \Lambda$
\[
\lambda |\xi|^2 \le a_{ij}\xi_i\xi_j \le \Lambda |\xi|^2
\qquad \text{for any } \xi \in \R^n.
\]

Suppose $w \in H^1(B_r(x_0))$ is a weak solution of
\[
(1)\qquad a_{ij}D_iD_j w = 0 \qquad \text{in } B_r(x_0).
\]

Then, for any $0<\rho\le r$
\[
\int_{B_\rho(x_0)} |Dw|^2 \le C\Bigl(\frac{\rho}{r}\Bigr)^n \int_{B_r(x_0)} |Dw|^2,
\]
and
\[
\int_{B_\rho(x_0)} \bigl|Dw-(Dw)_{x_0,\rho}\bigr|^2
\le C\Bigl(\frac{\rho}{r}\Bigr)^{n+2}\int_{B_r(x_0)} \bigl|Dw-(Dw)_{x_0,r}\bigr|^2,
\]
where $C$ is a positive constant depending only on $n$ and $\Lambda/\lambda$.

\textbf{Proof.} Note that if $w$ is a solution of (1) so is any of its derivatives. We may
apply Lemma 1.42 in Chapter 1 to $Dw$. \hfill $\square$
\end{lemma}

Now we compare any functions with harmonic functions.

\begin{corollary}\label{hanlin:weaksol-part1-cor-3-10}\dcref{Corollary 3.10}\group{ch03-weak-solutions-part-i}\source{han-lin-lecture-notes:page-0065}
\emph{Suppose $w$ is as in Lemma 2.2 and $u$ is an arbitrary $H^1$-function in $B_r(x_0)$. Then, for any $0<\rho\le r$}
\[
\int_{B_\rho(x_0)} |Du|^2 \le C\left\{\left(\frac{\rho}{r}\right)^n \int_{B_r(x_0)} |Du|^2 + \int_{B_r(x_0)} |D(u-w)|^2\right\},
\]
and
\[
\int_{B_\rho(x_0)} |Du-(Du)_{x_0,\rho}|^2 \le C\left\{\left(\frac{\rho}{r}\right)^{n+2} \int_{B_r(x_0)} |Du-(Du)_{x_0,r}|^2
+ \int_{B_r(x_0)} |D(u-w)|^2\right\},
\]
where $C$ is a positive constant depending only on $n$ and $\Lambda/\lambda$.
\end{corollary}
\begin{proof}
With $v=u-w$, we have for any $0<\rho\le r$
\[
\int_{B_\rho(x_0)} |Du|^2 \le 2\int_{B_\rho(x_0)} |Dw|^2 + 2\int_{B_\rho(x_0)} |Dv|^2
\]
\[
\le C\left(\frac{\rho}{r}\right)^n \int_{B_r(x_0)} |Dw|^2 + 2\int_{B_r(x_0)} |Dv|^2
\]
\[
\le C\left(\frac{\rho}{r}\right)^n \int_{B_r(x_0)} |Du|^2 + c\left[1+\left(\frac{\rho}{r}\right)^n\right]\int_{B_r(x_0)} |Dv|^2,
\]
and
\[
\int_{B_\rho(x_0)} |Du-(Du)_{x_0,\rho}|^2
\]
\[
\le 2\int_{B_\rho(x_0)} |Du-(Dw)_{x_0,\rho}|^2 + 2\int_{B_\rho(x_0)} |Dv|^2
\]
\[
\le 4\int_{B_\rho(x_0)} |Dw-(Dw)_{x_0,\rho}|^2 + 6\int_{B_\rho(x_0)} |Dv|^2
\]
\[
\le C\left(\frac{\rho}{r}\right)^{n+2} \int_{B_r(x_0)} |Dw-(Dw)_{x_0,r}|^2 + 6\int_{B_r(x_0)} |Dv|^2
\]
\[
\le C\left(\frac{\rho}{r}\right)^{n+2} \int_{B_r(x_0)} |Du-(Du)_{x_0,r}|^2 + c\left[1+\left(\frac{\rho}{r}\right)^{n+2}\right]\int_{B_r(x_0)} |Dv|^2.
\]
This finishes the proof.

The regularity of $u$ depends on how close $u$ is to $w$, the solution of the constant coefficient equation.

We now prove Theorem 3.8.
\end{proof}

\begin{proof}[Proof of Theorem 3.8]
We decompose $u$ into a sum $v+w$ where $w$ satisfies a homogeneous equation and $v$ has estimates in terms of nonhomogeneous terms.
For any $B_r(x_0)\subset B_1$, we write the equation in the following form
\[
\int_{B_1} a_{ij}(x_0)D_i u D_j\varphi
= \int_{B_1} f\varphi - cu\varphi + (a_{ij}(x_0)-a_{ij}(x))D_i u D_j\varphi.
\]
In $B_r(x_0)$, the Dirichlet problem
\[
\int_{B_r(x_0)} a_{ij}(x_0)D_i w D_j\varphi = 0 \qquad \text{for any } \varphi\in H^1_0(B_r(x_0))
\]
has a unique solution $w$ with $w-u\in H_0^1(B_r(x_0))$. Obviously, the function
$v=u-w\in H_0^1(B_r(x_0))$ satisfies the equation
\[
\int_{B_r(x_0)} a_{ij}(x_0)D_i vD_j\varphi
= \int_{B_r(x_0)} f\varphi-cuv+(a_{ij}(x_0)-a_{ij}(x))D_i uD_j\varphi
\]
for any $\varphi\in H_0^1(B_r(x_0))$.

By taking the test function $\varphi=v$, we obtain
\[
\int_{B_r(x_0)} |Dv|^2 \leq c\left\{\tau^2(r)\int_{B_r(x_0)} |Du|^2 +
\left(\int_{B_r(x_0)} |c|^n\right)^{\frac{2}{n}}\int_{B_r(x_0)} u^2
+\left(\int_{B_r(x_0)} |f|^{\frac{2n}{n+2}}\right)^{\frac{n+2}{n}}\right\},
\]
where we used the Sobolev inequality
\[
\left(\int_{B_r(x_0)} v^{\frac{2n}{n-2}}\right)^{\frac{n-2}{2n}}
\leq c(n)\left(\int_{B_r(x_0)} |Dv|^2\right)^{\frac12},
\]
for $v\in H_0^1(B_r(x_0))$. Therefore Corollary 3.10 implies for any $0<\rho\leq r$
\begin{equation}
\int_{B_\rho(x_0)} |Du|^2 \leq c\left\{\left[\left(\frac{\rho}{r}\right)^n+\tau^2(r)\right]\int_{B_r(x_0)} |Du|^2\right.
\left.+\left(\int_{B_r(x_0)} |c|^n\right)^{\frac{2}{n}}\int_{B_r(x_0)} u^2
+\left(\int_{B_r(x_0)} |f|^{\frac{2n}{n+2}}\right)^{\frac{n+2}{n}}\right\},
\tag{1}
\end{equation}
where $C$ is a positive constant depending only on $\lambda$ and $\Lambda$. By the Hölder inequality,
we have
\[
\left(\int_{B_r(x_0)} |f|^{\frac{2n}{n+2}}\right)^{\frac{n+2}{n}}
\leq \left(\int_{B_r(x_0)} |f|^q\right)^{\frac{2}{q}}r^{n-2+2\alpha},
\]
where $\alpha=2-n/q\in(0,1)$ if $n/2<q<n$. Hence (1) implies for any $B_r(x_0)\subset B_1$
and any $0<\rho\leq r$
\[
\int_{B_\rho(x_0)} |Du|^2 \leq C\left\{\left[\left(\frac{\rho}{r}\right)^n+\tau^2(r)\right]\int_{B_r(x_0)} |Du|^2+r^{n-2+2\alpha}\|f\|_{L^q(B_1)}^2
\right.
\]
\[
\left.+\left(\int_{B_r(x_0)} |c|^n\right)^{\frac{2}{n}}\int_{B_r(x_0)} u^2\right\}.
\]

\textit{Case 1.} $c\equiv 0$. We have for any $B_r(x_0)\subset B_1$ and for any $0<\rho\leq r$
\[
\int_{B_\rho(x_0)} |Du|^2 \leq C\left\{\left[\left(\frac{\rho}{r}\right)^n+\tau^2(r)\right]\int_{B_r(x_0)} |Du|^2+r^{n-2+2\alpha}\|f\|_{L^q(B_1)}^2\right\}.
\]

Now the result would follow if in the above inequality we could write $\rho^{n-2+2\alpha}$
instead of $r^{n-2+2\alpha}$. This is in fact true and is stated in Lemma 3.4. By Lemma
\end{proof}

\begin{proof}[Continued proof of Theorem 3.8]
3.4, there exists an $R_0 > 0$ such that, for any $x_0 \in B_{1/2}$ and any $0 < \rho < r \le R_0$,
\[
\int_{B_\rho(x_0)} |Du|^2 \le C\left\{\left(\frac{\rho}{r}\right)^{n-2+2\alpha}\int_{B_r(x_0)} |Du|^2 + r^{n-2+2\alpha}\|f\|_{L^q(B_1)}^2\right\}.
\]
In particular, taking $r = R_0$, we obtain for any $\rho < R_0$
\[
\int_{B_\rho(x_0)} |Du|^2 \le C\rho^{n-2+2\alpha}\left\{\int_{B_1} |Du|^2 + \|f\|_{L^q(B_1)}^2\right\}.
\]

\textit{Case 2.} General case. We have for any $B_r(x_0) \subset B_1$ and any $0 < \rho \le r$
\begin{equation}
\begin{aligned}
\int_{B_\rho(x_0)} |Du|^2 \le C\biggl\{&\left[\left(\frac{\rho}{r}\right)^n + \tau^2(r)\right]\int_{B_r(x_0)} |Du|^2 \\
&+ r^{n-2+2\alpha}\chi(F) + \int_{B_r(x_0)} u^2\biggr\}.
\end{aligned}
\tag{2}
\end{equation}
where $\chi(F) = \|f\|_{L^q(B_1)}^2$. We will prove for any $x_0 \in B_{1/2}$ and any $0 < \rho < r \le 1/2$
\begin{equation}
\begin{aligned}
\int_{B_\rho(x_0)} |Du|^2 \le C\biggl\{&\left[\left(\frac{\rho}{r}\right)^n + \tau^2(r)\right]\int_{B_r(x_0)} |Du|^2 \\
&+ r^{n-2+2\alpha}\left[\chi(F) + \int_{B_1} u^2 + \int_{B_1} |Du|^2\right]\biggr\}.
\end{aligned}
\tag{3}
\end{equation}

We need a bootstrap argument. First by Lemma 3.3, there exists an $R_1 \in (1/2,1)$
such that for any $x_0 \in B_{R_1}$ and any $0 < r \le 1 - R_1$
\begin{equation}
\int_{B_r(x_0)} u^2 \le Cr^{\delta_1}\left\{\int_{B_1} |Du|^2 + \int_{B_1} u^2\right\},
\tag{4}
\end{equation}
where $\delta_1 = 2$ if $n > 2$ and $\delta_1$ is arbitrary in $(0,2)$ if $n = 2$. This, with (2), yields
\begin{equation}
\begin{aligned}
\int_{B_\rho(x_0)} |Du|^2 \le C\biggl\{&\left[\left(\frac{\rho}{r}\right)^n + \tau^2(r)\right]\int_{B_r(x_0)} |Du|^2 \\
&+ r^{n-2+2\alpha}\chi(F) + r^{\delta_1}\|u\|_{H^1(B_1)}^2\biggr\}.
\end{aligned}
\end{equation}

Then (3) holds in the following cases: (i) $n = 2$, by choosing $\delta_1 = 2\alpha$; (ii) $n > 2$
while $n - 2 + 2\alpha \le 2$, by choosing $\delta_1 = 2$. For $n > 2$ and $n - 2 + 2\alpha > 2$, we have
\[
\int_{B_\rho(x_0)} |Du|^2 \leq c\left\{\left[\left(\frac{\rho}{r}\right)^n + \tau^2(r)\right]\int_{B_r(x_0)} |Du|^2 + r^2\left[\chi(F) + \|u\|_{H^1(B_1)}^2\right]\right\}.
\]
Lemma 3.4 again yields for any $x_0 \in B_{R_1}$ and any $0 < r \le 1 - R_1$
\[
\int_{B_r(x_0)} |Du|^2 \leq Cr^2\left\{\chi(F) + \|u\|_{H^1(B_1)}^2\right\}.
\]
Hence by Lemma 3.3, there exists an $R_2 \in (1/2,R_1)$ such that for any $x_0 \in B_{R_2}$
and any $0 < r \le R_1 - R_2$
\begin{equation}
\int_{B_r(x_0)} u^2 \leq Cr^{\delta_2}\left\{\chi(F) + \|u\|_{H^1(B_1)}^2\right\},
\tag{5}
\end{equation}
where $\delta_2 = 4$ if $n > 4$ and $\delta_2$ is arbitrary in $(2,n)$ if $n = 3$ or $4$. Notice (5) is an improvement compared with (4). Substitute (5) in (2) and continue the process. After finitely many steps, we obtain (3). $\square$
\end{proof}

\subsection*{3.3. Hölder Continuity of Gradients}

In this section, we prove the Hölder regularity for gradients of solutions. We follow the same idea used to prove Theorem 3.8.

Suppose $a_{ij} \in L^\infty(B_1)$ is uniformly elliptic in $B_1$, i.e.,
\[
\lambda |\xi|^2 \leq a_{ij}(x)\xi_i\xi_j \leq \Lambda |\xi|^2 \qquad \text{for any } x \in B_1,\, \xi \in \mathbb{R}^n.
\]

We assume that $u \in H^1(B_1)$ satisfies
\[
(*) \qquad \int_{B_1} a_{ij}D_i u D_j \varphi + cu\varphi = \int_{B_1} f\varphi \qquad \text{for any } \varphi \in H^1_0(B_1).
\]

The main theorem in this section is the following Hölder estimate for gradients.

\begin{theorem}\label{hanlin:weaksol-part1-grad-holder}\dcref{Theorem 3.11}\group{ch03-weak-solutions-part-i}\source{han-lin-lecture-notes:page-0068,han-lin-lecture-notes:page-0069,han-lin-lecture-notes:page-0070,han-lin-lecture-notes:page-0071}
Let $u \in H^1(B_1)$ solve (*). Assume $a_{ij} \in C^\alpha(B_1)$, $c \in L^q(B_1)$ and $f \in L^q(B_1)$ for some $q > n$ and $\alpha = 1 - n/q \in (0,1)$. Then $Du \in C^\alpha(B_1)$. Moreover, there exists an $R_0 \in (0,1)$ such that, for any $x \in B_{1/2}$ and $r \leq R_0$,
\[
\int_{B_r(x)} \lvert Du - (Du)_{x,r}\rvert^2 \leq Cr^{n+2\alpha}\left\{\lVert f\rVert_{L^q(B_1)}^2 + \lVert u\rVert_{H^1(B_1)}^2\right\},
\]
where $R_0$ and $C$ are positive constants depending only on $\lambda$, $\lvert a_{ij}\rvert_{C^\alpha}$ and $\lvert c\rvert_{L^q}$.
\end{theorem}
\begin{proof}
As in the proof of Theorem 3.8, we decompose $u$ into a sum $v+w$ where $w$ satisfies a homogeneous equation and $v$ has estimates in terms of nonhomogeneous terms.

For any $B_r(x_0) \subset B_1$, we write the equation in the following form
\[
\int_{B_1} a_{ij}(x_0)D_i u D_j \varphi
= \int_{B_1} f\varphi - cu\varphi + \bigl(a_{ij}(x_0) - a_{ij}(x)\bigr)D_i u D_j \varphi.
\]

In $B_r(x_0)$, the Dirichlet problem
\[
\int_{B_r(x_0)} a_{ij}(x_0)D_i w D_j \varphi = 0 \qquad \text{for any } \varphi \in H^1_0(B_r(x_0)),
\]
has a unique solution $w$ with $w - u \in H^1_0(B_r(x_0))$. Obviously, the function $v = u - w \in H^1_0(B_r(x_0))$ satisfies the equation
\[
\int_{B_r(x_0)} a_{ij}(x_0)D_i v D_j \varphi
= \int_{B_r(x_0)} f\varphi - cu\varphi + \bigl(a_{ij}(x_0) - a_{ij}(x)\bigr)D_i u D_j \varphi
\]
for any $\varphi \in H^1_0(B_r(x_0))$.

By taking the test function $\varphi = v$, we obtain
\[
\int_{B_r(x_0)} \lvert Dv\rvert^2 \leq C\left\{\tau^2(r)\int_{B_r(x_0)} \lvert Du\rvert^2 + \left(\int_{B_r(x_0)} \lvert c\rvert^n\right)^{\frac{2}{n}}\int_{B_r(x_0)} u^2 + \left(\int_{B_r(x_0)} \lvert f\rvert^{\frac{2n}{n+2}}\right)^{\frac{n+2}{n}}\right\}.
\]
\end{proof}

\begin{proof}[Continued proof of Theorem 3.11]
Therefore Corollary 3.10 implies for any $0<\rho\le r$
\begin{equation}
\int_{B_\rho(x_0)} \lvert Du\rvert^2 \le C\left\{\left[\left(\frac{\rho}{r}\right)^n+\tau^2(r)\right]\int_{B_r(x_0)} \lvert Du\rvert^2
+ \left(\int_{B_r(x_0)} \lvert c\rvert^n\right)^{\frac{2}{n}}\int_{B_r(x_0)} u^2 + \left(\int_{B_r(x_0)} \lvert f\rvert^{\frac{2n}{n+2}}\right)^{\frac{n+2}{n}}\right\},
\tag{1}
\end{equation}
and
\begin{equation}
\int_{B_\rho(x_0)} \lvert Du-(Du)_{x_0,\rho}\rvert^2
\le C\left\{\left(\frac{\rho}{r}\right)^{n+2}\int_{B_r(x_0)} \lvert Du-(Du)_{x_0,r}\rvert^2+\tau^2(r)\int_{B_r(x_0)} \lvert Du\rvert^2
+ \left(\int_{B_r(x_0)} \lvert c\rvert^n\right)^{\frac{2}{n}}\int_{B_r(x_0)} u^2 + \left(\int_{B_r(x_0)} \lvert f\rvert^{\frac{2n}{n+2}}\right)^{\frac{n+2}{n}}\right\},
\tag{2}
\end{equation}
where $C$ is a positive constant depending only on $\lambda$ and $A$.

By the Hölder inequality, we have for any $B_r(x_0)\subset B_1$
\begin{equation}
\left(\int_{B_r(x_0)} \lvert f\rvert^{\frac{2n}{n+2}}\right)^{\frac{n+2}{n}}
\le \left(\int_{B_r(x_0)} \lvert f\rvert^q\right)^{\frac{2}{q}} r^{n+2\alpha},
\end{equation}
and
\begin{equation}
\left(\int_{B_r(x_0)} \lvert c\rvert^n\right)^{\frac{2}{n}} \le r^{2\alpha}\left(\int_{B_r(x_0)} \lvert c\rvert^q\right)^{\frac{2}{q}},
\end{equation}
with $\alpha=1-\frac{n}{q}$.

Case 1. $a_{ij}\equiv \mathrm{const.}$, $c\equiv 0$. In this case, $\tau(r)\equiv 0$. By (2), there holds for any $B_r(x_0)\subset B_1$ and $0<\rho\le r$
\begin{equation}
\int_{B_\rho(x_0)} \lvert Du-(Du)_{x_0,\rho}\rvert^2
\le C\left\{\left(\frac{\rho}{r}\right)^{n+2}\int_{B_r(x_0)} \lvert Du-(Du)_{x_0,r}\rvert^2+r^{n+2\alpha}\|f\|_{L^q(B_1)}^2\right\}.
\end{equation}
By Lemma 3.4, we replace $r^{n+2\alpha}$ by $\rho^{n+2\alpha}$ to get the result.

Case 2. $c\equiv 0$. By (1) and (2), we have for any $B_r(x_0)\subset B_1$ and any $\rho<r$
\begin{equation}
\int_{B_\rho(x_0)} \lvert Du\rvert^2 \le C\left\{\left[\left(\frac{\rho}{r}\right)^n+r^{2\alpha}\right]\int_{B_r(x_0)} \lvert Du\rvert^2+r^{n+2\alpha}\|f\|_{L^q(B_1)}^2\right\},
\tag{3}
\end{equation}
and
\begin{equation}
\int_{B_\rho(x_0)} \lvert Du-(Du)_{x_0,\rho}\rvert^2 \le C\left\{\left(\frac{\rho}{r}\right)^{n+2}\int_{B_r(x_0)} \lvert Du-(Du)_{x_0,r}\rvert^2
+ r^{2\alpha}\int_{B_r(x_0)} \lvert Du\rvert^2 + r^{n+2\alpha}\|f\|_{L^q(B_1)}^2\right\}.
\tag{4}
\end{equation}

We need to estimate the integral
\[
\int_{B_r(x_0)} |Du|^2.
\]

Write $\chi(F) = \|f\|_{L^q(B_1)}^2$.
Take small $\delta > 0$. Then (3) implies
\[
\int_{B_\rho(x_0)} |Du|^2 \le C \left\{ \left[\left(\frac{\rho}{r}\right)^n + r^{2\alpha}\right] \int_{B_r(x_0)} |Du|^2 + r^{n-2\delta} \chi(F) \right\}.
\]

Hence Lemma 3.4 implies the existence of an $R_1 \in (3/4,1)$ with $r_1 = 1 - R_1$ such that, for any $x_0 \in B_{R_1}$ and any $0 < r \le r_1$,
\[
(5)\qquad \int_{B_r(x_0)} |Du|^2 \le C r^{n-2\delta} \left\{ \chi(F) + \|Du\|_{L^2(B_1)}^2 \right\}.
\]

Therefore, by substituting (5) in (4) we obtain for any $0 < \rho < r \le r_1$
\[
\int_{B_\rho(x_0)} |Du - (Du)_{x_0,\rho}|^2 \le C \left\{ \left(\frac{\rho}{r}\right)^{n+2} \int_{B_r(x_0)} |Du - (Du)_{x_0,r}|^2
+ r^{n+2\alpha-2\delta} \left[ \chi(F) + \|Du\|_{L^2(B_1)}^2 \right] \right\}.
\]

By Lemma 3.4 again, we have for any $x_0 \in B_{R_1}$ and any $0 < \rho < r \le r_1$
\[
\int_{B_\rho(x_0)} |Du - (Du)_{x_0,\rho}|^2 \le C \left\{ \left(\frac{\rho}{r}\right)^{n+2\alpha-2\delta} \int_{B_r(x_0)} |Du - (Du)_{x_0,r}|^2
+ \rho^{n+2\alpha-2\delta} \left[\chi(F) + \|Du\|_{L^2(B_1)}^2 \right] \right\}.
\]

With $r = r_1$ this implies that for any $x_0 \in B_{R_1}$ and any $0 < r \le r_1$
\[
\int_{B_r(x_0)} |Du - (Du)_{x_0,r}|^2 \le C r^{n+2\alpha-2\delta} \left\{ \chi(F) + \|Du\|_{L^2(B_1)}^2 \right\}.
\]

Hence $Du \in C_{\mathrm{loc}}^{\alpha-\delta}$ for any $\delta > 0$ small. In particular, $Du \in L^\infty_{\mathrm{loc}}$ and
\[
(6)\qquad \sup_{B_{3/4}} |Du|^2 \le C \left\{ \chi(F) + \|Du\|_{L^2(B_1)}^2 \right\}.
\]

Combining (4) and (6), we have for any $x_0 \in B_{1/2}$ and $0 < \rho < r \le r_1$
\[
\int_{B_\rho(x_0)} |Du - (Du)_{x_0,\rho}|^2 \le C \left\{ \left(\frac{\rho}{r}\right)^{n+2} \int_{B_r(x_0)} |Du - (Du)_{x_0,r}|^2
+ r^{n+2\alpha} \left[\chi(F) + \|Du\|_{L^2(B_1)}^2 \right] \right\}.
\]

By Lemma 3.4 again, this implies
\[
\int_{B_\rho(x_0)} |Du - (Du)_{x_0,\rho}|^2 \le C \left\{ \left(\frac{\rho}{r}\right)^{n+2\alpha} \int_{B_r(x_0)} |Du - (Du)_{x_0,r}|^2
+ \rho^{n+2\alpha} \left[\chi(F) + \|Du\|_{L^2(B_1)}^2 \right] \right\}.
\]

\[
\text{By choosing } r = r_1, \text{ we have for any } x_0 \in B_{\frac12} \text{ and } r \le r_1
\]
\[
\int_{B_r(x_0)} \lvert Du - (Du)_{x_0,r} \rvert^2 \le c r^{n+2\alpha} \left\{ \chi(F) + \lVert Du \rVert_{L^2(B_1)}^2 \right\}.
\]

\[
\textit{Case 3. The general case. By (1) and (2), we have for any } B_r(x_0) \subset B_1 \text{ and }
\]
\[
\rho < r
\]
\[
\tag{7}
\int_{B_\rho(x_0)} \lvert Du \rvert^2 \le C \left\{ \left[ \left( \frac{\rho}{r} \right)^n + r^{2\alpha} \right] \int_{B_r(x_0)} \lvert Du \rvert^2
+ \int_{B_r(x_0)} u^2 + r^{n+2\alpha}\chi(F) \right\},
\]
and
\[
\tag{8}
\int_{B_\rho(x_0)} \lvert Du - (Du)_{x_0,\rho} \rvert^2 \le C \left\{ \left( \frac{\rho}{r} \right)^{n+2} \int_{B_r(x_0)} \lvert Du - (Du)_{x_0,r} \rvert^2 \right.
\]
\[
\left. +\ r^{2\alpha} \left[ \int_{B_r(x_0)} u^2 + \int_{B_r(x_0)} \lvert Du \rvert^2 \right] + r^{n+2\alpha}\chi(F) \right\},
\]
\[
\text{where } \chi(F) = \lVert f \rVert_{L^q(B_1)}^2. \text{ In (7), we replace } r^{n+2\alpha} \text{ by } r^n. \text{ As in the proof of}
\]
\[
\text{Theorem 3.8, we show that for any small } \delta > 0 \text{ there exists an } R_1 \in (3/4,1) \text{ such}
\]
\[
\text{that for any } x \in B_{R_1} \text{ and } r < 1 - R_1
\]
\[
\tag{9}
\int_{B_r(x_0)} \lvert Du \rvert^2 \le C r^{n-2\delta} \left\{ \chi(F) + \lVert u \rVert_{H_1^1(B)}^2 \right\}.
\]
\[
\text{By Lemma 3.3, we also get}
\]
\[
\tag{10}
\int_{B_r(x_0)} u^2 \le C r^{n-2\delta} \left\{ \chi(F) + \lVert u \rVert_{H_1^1(B)}^2 \right\}.
\]
\[
\text{Write}
\]
\[
\chi(F,u) = \lVert f \rVert_{L^q}^2 + \lVert u \rVert_{H^1}^2.
\]
\[
\text{Then, (8), (9) and (10) imply}
\]
\[
\int_{B_\rho(x_0)} \lvert Du - (Du)_{x_0,\rho} \rvert^2
\]
\[
\le C \left\{ \left( \frac{\rho}{r} \right)^{n+2} \int_{B_r(x_0)} \lvert Du - (Du)_{x_0,r} \rvert^2 + r^{n+2\alpha-2\delta}\chi(F,u) \right\}.
\]
\[
\text{Hence, Lemma 3.4 and Theorem 3.1 imply that } Du \in C_{\mathrm{loc}}^{\alpha-\delta} \text{ for small } \delta < \alpha. \text{ In}
\]
\[
\text{particular, we have } u \in C_{\mathrm{loc}}^1 \text{ with}
\]
\[
\tag{11}
\sup_{B_{3/4}} \lvert u \rvert^2 + \sup_{B_{3/4}} \lvert Du \rvert^2 \le C \chi(F,u).
\]
\[
\text{Now (8) and (11) imply}
\]
\[
\int_{B_\rho(x_0)} \lvert Du - (Du)_{x_0,\rho} \rvert^2
\]
\[
\le C \left\{ \left( \frac{\rho}{r} \right)^{n+2} \int_{B_r(x_0)} \lvert Du - (Du)_{x_0,r} \rvert^2 + r^{n+2\alpha}\chi(F,u) \right\}.
\]
\[
\text{This finishes the proof.}
\]
\end{proof}

It is natural to ask whether $f\in L^\infty(B_1)$, with appropriate assumptions on $a_{ij}$
and $c$, implies $Du\in C^1$. Consider
\[
\int_{B_1} D_i u\, D_i\varphi = \int_{B_1} f\varphi \qquad \text{for any } \varphi \in H^1_0(B_1).
\]
There exists an example showing that $f\in C$ and $u\in C^{1,\alpha}_{loc}$ for any $\alpha\in (0,1)$ while
$D^2u\notin C$.

\begin{example}\label{hanlin:weaksol-part1-example-3-12}\dcref{Example 3.12}\group{ch03-weak-solutions-part-i}\source{han-lin-lecture-notes:page-0072}
In $B_R \subset \mathbb{R}^n$, with $R<1$, consider
\[
\Delta u = \frac{x_2^2 - x_1^2}{2|x|^2}\left\{\frac{n+2}{(-\log |x|)^{1/2}} + \frac{1}{2(-\log |x|)^{3/2}}\right\},
\]
where the right hand side is continuous in $\overline{B}_R$ if we set it equal to zero at the origin.
The function $u(x) = (x_1^2 - x_2^2)(-\log |x|)^{1/2} \in C(\overline{B}_R)\cap C^\infty(\overline{B}_R\setminus\{0\})$ satisfies
the above equation in $B_R\setminus\{0\}$ and the boundary condition $u = \sqrt{-\log R}(x_1^2 - x_2^2)$ on
$\partial B_R$. Note that $u$ is not a classical solution of the problem since $\lim_{|x|\to 0}D_{11}u=\infty$,
and therefore $u$ is not in $C^2(B_R)$. In fact, the problem has no classical solutions (al-
though it has a weak solution). Assume on the contrary that a classical solution $v$
exists. Then the function $w = u-v$ is harmonic and bounded in $B_R\setminus\{0\}$. By The-
orem 1.28 in Chapter 1, on removable singularities for harmonic functions, $w$ may
be redefined at the origin so that $\Delta w = 0$ in $B_R$ and therefore belongs to $C^2(B_R)$.
In particular, the (finite) limit $\lim_{|x|\to 0} D_{11}u$ exists, which is a contradiction.
\end{example}
